\documentclass[11pt,a4paper]{article}
\usepackage[margin=20mm]{geometry}
\usepackage{fontspec}
\setmainfont{DejaVu Serif}
\setsansfont{DejaVu Sans}
\setmonofont{DejaVu Sans Mono}[Scale=0.85]
\usepackage{polyglossia}
\setdefaultlanguage{russian}
\usepackage{graphicx}
\usepackage{hyperref}
\usepackage{enumitem}
\usepackage{fancyvrb}
\hypersetup{hidelinks}
\setlength{\parindent}{0pt}
\setlength{\parskip}{4pt}
\usepackage{titlesec}
\titleformat{\section}{\large\bfseries}{}{0pt}{}
\titlespacing*{\section}{0pt}{10pt}{4pt}
\sloppy
\setlist{nosep,leftmargin=*}
\emergencystretch=3em
\newcommand{\shot}[2]{\begin{center}\includegraphics[width=\linewidth,height=.34\textheight,keepaspectratio]{#1}\par\small #2\end{center}}
\begin{document}
\begin{center}
{\small Университет ИТМО · Фронт-энд разработка · 2026}\par
{\LARGE Домашняя работа № 2}\par
{\large Доступность HTML-интерфейса}\par
Светлов Ярослав, группа J3211
\end{center}
\section*{Цель}
Улучшить доступность самостоятельной копии ЛР2 и проверить её с помощью Lighthouse и Firefox Accessibility Inspector.
\section*{Изменения}
\begin{itemize}
\item Ссылка перехода к содержимому, видимый фокус, перенос фокуса на заголовок при смене страницы.
\item Видимые подписи полей, заголовки страниц, aria-current для навигации и aria-busy при загрузке.
\item Ошибки и уведомления используют alert/status. Кнопки сохранения имеют разные доступные имена.
\item У таблиц есть caption и заголовки столбцов. Прокручиваемая область получает Tab-фокус только при переполнении.
\item График продублирован численной таблицей; отключено взаимодействие, доступное только мышью.
\item Усилен контраст ссылок и кнопок. Дата вводится в подписанном формате ГГГГ-ММ-ДД с проверкой: встроенные части календаря Firefox вызывали замечания инспектора.
\end{itemize}

\section*{Результаты аудита}
Lighthouse проверил шесть экранов: вход, регистрацию, кабинет, поиск, эксперимент и модели. Все получили 100 баллов в категории доступности; результаты сохранены в \texttt{checks/lighthouse.json}.

В Firefox 148 выполнен аудит ALL через Accessibility Inspector на тех же экранах. После исправлений нет FAIL. Осталось предупреждение о видимости фокуса уже сфокусированного поля Email; переход Tab проверен отдельно, контрастный outline задан в CSS. Результаты: \texttt{checks/firefox-accessibility.json}.

Проверены клавиатурная навигация, вход, пустая выдача и сброс. Девять API-тестов пройдены.
\section*{Запуск и проверка}
Требуется Node.js 22.12 или новее. Команды выполняются из папки этой работы.
\begin{verbatim}
npm ci
npm start
\end{verbatim}
Открыть \url{http://localhost:3002}. Автотесты: \texttt{npm test}.
Учебные аккаунты: \texttt{admin@nexus.ai / password} и \texttt{vision@nexus.ai / cv12345}. Данные сохраняются в \texttt{db.local.json}, исходная база \texttt{db.json} не меняется.\section*{Вывод и ограничения}
Доступность основных экранов улучшена и проверена двумя инструментами. Полное тестирование со скринридером не проводилось. Автоматические оценки не доказывают полного соответствия WCAG.

\clearpage
\section*{Скриншоты}
\shot{checks/firefox-inspector.png}{Рисунок 1. Аудит в Firefox Accessibility Inspector.}
\shot{checks/search.png}{Рисунок 2. Поиск после исправлений доступности.}

\end{document}
